\chapter*{الملخص}

يقدم هذا المشروع دراسة تحليلية وتطبيقية شاملة للنمذجة الرياضية وتصميم استراتيجيات التحكم المتقدمة للروبوت الصناعي ذي المفاصل الأفقية \LR{(SCARA)}.     يهدف البحث إلى تحقيق تتبع مسار دقيق وسريع يواكب التغيرات التشغيلية والاضطرابات الخارجية.

يبدأ المشروع بدراسة النمذجة الحركية \LR{(Kinematic Modeling)} للروبوت، حيث تم استنتاج معادلات الحركة المباشرة \LR{(Forward Kinematics)} لتحديد موضع وتوجيه المؤثر النهائي في الفراغ التشغيلي بناءً على إحداثيات المفاصل. وفي الجانب التحليلي الأعمق، تم اشتقاق النموذج الديناميكي بدقة عالية باستخدام معادلات لاغرانج-أويلر، حيث يتميز هذا العمل بإدراج تغير كتلة الحمل \LR{$m_L$} بشكل صريح ومباشر كمتغير أساسي مدمج داخل بنية النموذج ومُقَدِّر الديناميكا اللحظي. هذا الإدراج المباشر مكَّن المنظومة من تحديث حسابات مصفوفة العطالة وشعاع الجاذبية آنياً بناءً على القيمة اللحظية للحمل، مما يضمن محاكاة فيزيائية دقيقة تواكب الحمولات الحقيقية للروبوت.

وبناءً على هذا النموذج الرياضي المطور، تم تصميم ومحاكاة ثلاث استراتيجيات تحكم رئيسية وتقييم أدائها؛ حيث طُبِّق متحكم تناسبي-تكاملي-تفاضلي \LR{(PID)} كمعيار مرجعي لبيان عيوب تراكم الأخطاء وتأثير الجاذبية، تلاه تصميم متحكم بالانزلاق \LR{(SMC)} لتوجيه الذراع بنعومة وفصل ديناميكا المحاور. ولتوفير أقصى درجات المتانة ضد أخطاء النمذجة المتبقية أو الاحتكاك غير المنمذج، تم تطوير متحكم عزم محسوب تكيفي متقدم يعتمد على مصفوفة الانحدار \LR{(Regressor-Based ACTC)}، مدعوماً بقانون تحديث يشتمل على عامل التسريب \LR{)$\sigma$-modification)} لضمان استقرار المعاملات المقدرة ومنع انجرافها رياضياً.

تمت محاكاة النظام بالكامل في بيئة \LR{MATLAB/Simulink}، حيث أثبتت النتائج التطبيقية كفاءة استراتيجيات التحكم المعتمدة على النموذج \LR{(Model-Based)} وقدرتها الفائقة على تحقيق استجابة ديناميكية ناعمة ومستقرة تماماً بدون تجاوز \LR{(Overshoot)}، مع تصفير خطأ الحالة المستقرة للمفصل الانسحابي والمفاصل الدورانية تحت مختلف ظروف التشغيل والحمولات المتغيرة.